\section{Discussion: make readiness earn its cost}
\subsection{What the intelligence-transition question becomes}
The motivating question is deliberately large. The answer is useful only after becoming small enough to test: whose bottleneck, what output, which deadline, what proof, which alternative and at what total cost? Sector scale alone supplies none of those quantities. The audited macro envelope cannot decide whether a single solver should be prepared today.

The experiment demonstrates a practical distinction between \emph{capacity} and \emph{readiness}. A solver available from a library can solve the system; a retained factorization and known verification recipe can respond sooner to a repeated request. That change in readiness has an observed latency consequence. Yet the mechanism is available to a conventional engineer. Comparing only against repeated setup would overstate the contribution of foresight and of the certification wrapper.

The same lesson appears at the policy layer. Information can matter near a build threshold, but it is not valuable merely because it sharpens a probability. In the uninformative regime the policy correctly declines acquisition. In the nominal regime its uncertain gain over a fixed rule does not survive the stronger claim that the proposed procedure is uniquely better: a conventional compiled rule makes identical decisions. The positive candidate mechanism is therefore ordinary preparation plus decision-relevant information, deployed efficiently. The label \TI{} is a motivating discipline, not an additional source of measured value.

\subsection{Waiting, competition and distribution change}
Waiting has two different meanings in this study. Mandatory waiting followed by construction can miss a short deadline or buy an expensive capability too late. Adaptive waiting can preserve an option to use an alternative, build only when requests justify the cost, or decline unproductive work. A credible alternative must retain those choices. The better performance of adaptive policies over the wait-then-build ablation is not a finding that patience is irrational.

Imitation removes exclusivity before it removes usefulness. When competing capability becomes cheaply accessible, a beneficiary can gain even as the original builder loses an appropriable advantage. Fast imitation therefore improves available service routes in the model; it does not retroactively erase already delivered utility. Price compression is represented only through an explicitly labeled fixed-trajectory value sensitivity. The experiment does not estimate a market-clearing price, demand elasticity, monopoly duration or investor return.

The reversed-signal regime gives a more direct warning. A model can continue to produce coherent posterior decisions after its likelihood assumptions become wrong. In this case the nominally informative intervention reduces modeled net value. The response is not to reinterpret the loss as hidden success. It is to monitor the validity of the information model, retain an alternative policy, and require new evidence before committing to a materially changed plan.

\subsection{The strongest useful method actually found}
For the tested numerical workloads, use conventional factor reuse with per-output checking. Retain a source-bound reconstruction recipe when auditability and later restoration are needed. Preserve a calibrated delivery distribution that distinguishes already-paid preparation from a cold start. Use the simplest competent policy that implements the justified build, rent or wait decision, and charge its model preparation and execution. Require fresh capability evidence and separate authority when invoking a controlled action.

This recommendation is bounded to the measured instrument and finite model. It does not assert that compiled finite rules will dominate general agents on open-ended research. It does show what an agent must beat here. Repeating a deliberation that has already been compiled adds cost without changing the useful action. Wrapping a millisecond-scale kernel in seconds of lifecycle bookkeeping is also a poor default unless that bookkeeping buys a necessary assurance. Those are actionable engineering findings even when incremental commercial value is unknown.

\section{Limitations and the next decisive evidence}
The four matrices and manufactured inputs are a small benchmark, not a representative sample of scientific missions. Shared-host timing varies; there is no cross-machine or large-scale generalization result. Specialized Cholesky, tuned iterative methods and alternative orderings were not exhaustively compared. The bounded CG diagnostic is not a substitute for such a study. The instrument's known-solution error check is available in manufactured tests but generally unavailable for a novel physical problem.

The local forecasts resolve after issuance, but their deadlines are relative to known computational starts and often distinguish very easy from very hard conditions. The pooled constant is a limited reference. This is prospective local execution evidence, not a historical language-model backtest or a forecast of industrial adoption. Public matrices may already be known to language models, and assistant-assisted task design is not an independently blinded human selection process.

Demand, adoption, imitation, delay, amortization and mission-credit values are assumed. Model-conditional intervals do not cover uncertainty about those assumptions. The finite policy family cannot stand in for every competent organization. Native role separation uses a public shared test identity, not authentic independent actors. Fleet execution, physical instrumentation, fresh mainnet state, live funding and transaction behavior are outside the demonstrated boundary. Contract source continuity does not prove deployed-state continuity.

The cost ledger is auditable where measurements exist and explicitly incomplete elsewhere. A positive operational difference would still need to pay authoring, acquisition, engineering, review, maintenance, adoption and opportunity costs. No recruited beneficiary, accepted public pilot or payment receipt closes that gap. The original seven- and thirty-day adoption targets remain unmeasured; they are not generated as synthetic successes, and the paper does not score a future outcome before it occurs.

The next decisive study would recruit a numerical-research beneficiary, freeze a real workload and acceptance predicate, and prospectively assign comparable tasks to the readiness procedure and a competent conventional workflow. Both sides would receive the same tools and information budget; engineering and review effort would be independently recorded. The outcome would be accepted investigations delivered by a real deadline, with maintenance and full labor accounted for. That extension is a proposal, not part of the reported evidence, and no outreach or deployment is implied.

\section{Conclusion}
Could the intelligence transition be an extraordinary opportunity? It could. This study does not measure its total size or forecast which sector wins first. It identifies a more demanding unit of progress: an accepted capability that is ready in time and worth its full cost relative to a credible response.

The numerical artifact works. Its local delivery forecasts are testable and show skill. Its restoration and authority boundaries can be inspected and exercised. But the comparative study does not establish a unique full-cost advantage for constructive foresight: conventional reuse explains the numerical gain, and an equally informed compiled policy matches the decisions with less overhead. That boundary is the result, not a defect to hide.

\begin{claimbox}
\textbf{Build the capability. Prove the delivery. Make readiness earn its cost.}\\[3pt]
Useful foresight ends in a qualified action---including the decision to wait---whose incremental value survives a serious alternative.
\end{claimbox}
